% Part: second-order-logic
% Chapter: metatheory
% Section: undecidability-and-axiomatizability

\documentclass[../../../include/open-logic-section]{subfiles}

\begin{document}

\olfileid{sol}{met}{nax}

\olsection{द्वितीय-क्रम तर्कशास्त्र स्वयंसिद्धकीकरणीय नाही}


\begin{thm}
\ollabel{thm:sol-undecidable}
द्वितीय-क्रम तर्कशास्त्र अनिर्णेय आहे.
\end{thm}

\begin{proof}
प्रथम-क्रम !!{sentence} प्रथम-क्रम तर्कशास्त्रात वैध असेल,
जर आणि फक्त जर ते द्वितीय-क्रम तर्कशास्त्रात वैध असेल.
प्रथम-क्रम तर्कशास्त्र अनिर्णेय आहे.
\end{proof}

\begin{thm}
\ollabel{cor:sol-not-axiomatizable}
द्वितीय-क्रम तर्कशास्त्राच्या मानक चिन्हार्थमीमांसेसाठी
निर्दोष, संपूर्ण आणि प्रभावी !!{derivation} प्रणाली अस्तित्वात नाही.
\end{thm}

\begin{proof}
%READERNOTE{SOL6-A095: गैरस्वयंसिद्धकीकरणाचा निष्कर्ष प्रभावी सिद्धता प्रणालींविषयी आहे. Reduction मध्ये प्रथम क्रमाची अंकगणितीय वाक्ये वापरली असून नव्या चलांची निवड capture टाळणारी आहे.}
अंकगणिताच्या प्रथम-क्रम भाषेतील !!a{sentence}~$!A$ असो.
$\Sat{N}{!A}$ जर आणि फक्त जर $\Th{PA^2} \Entails !A$.
$!P$ हा~$\Th{PA^2}$ च्या नऊ स्वयंसिद्धकांचा संयोग असो.
$\Th{PA^2} \Entails !A$ जर आणि फक्त जर
$\Entails !P \lif !A$; म्हणजे प्रत्येक रचनेसाठी
$\Sat{M}{!P \lif !A}$. आता पुढील !!{sentence} पाहू:
$\lforall[z][\lforall[u][\lforall[u'][\lforall[u''][\lforall[L][(!P' \lif !A')]]]]]$.
हे $\Obj{0}$ च्या जागी $z$, $\prime$ च्या जागी
एकस्थानी फलनचल $u$, $+$ आणि $\times$ यांच्या जागी
अनुक्रमे द्विस्थानी फलनचले $u'$ आणि $u''$,
आणि $<$ च्या जागी द्विस्थानी संबंधचल~$L$ ठेवून
नंतर सार्वत्रिक संख्यापनाने मिळते. येथे नव्याने वापरलेले
सर्व चले आधीच्या सूत्रांत नसलेले निवडतो; आवश्यक असल्यास
आधी बद्ध चलांची नावे बदलतो, म्हणून capture होत नाही. हे शुद्ध द्वितीय-क्रम
तर्कशास्त्रातील वैध वाक्य असेल, जर आणि फक्त जर मूळ
वाक्य वैध असेल, जर आणि फक्त जर
$\Th{PA^2} \Entails !A$, जर आणि फक्त जर
$\Sat{N}{!A}$. म्हणून द्वितीय-क्रम तर्कशास्त्रासाठी
निर्दोष, संपूर्ण आणि प्रभावी सिद्धता-प्रणाली असती, तर
प्रथम-क्रम अंकगणितीय वाक्यांचे आणि त्यांच्या वरील निर्मित
वाक्यांसाठीच्या सिद्धतांचे समांतर प्रभावी प्रगणन वापरून
$\Struct{N}$ मध्ये सत्य असलेल्या !!{sentence}s चे
संगणनक्षम प्रगणन $f\colon \Nat \to \Sent[L_A]$
परिभाषित करता आले असते. हे फलन काही प्रथम-क्रम
सूत्र~$!B_f(x, y)$ द्वारे~$\Th{Q}$ मध्ये निरूपित
करता आले असते. मग !!{formula}~$\lexists[x][!B_f(x, y)]$
हे~$\Struct{N}$ मधील सत्य प्रथम-क्रम !!{sentence}s चा संच
परिभाषित करील. हा टार्स्कीच्या प्रमेयाशी विरोध आहे.
\end{proof}



\end{document}
